% !TEX root = ../../mixed-sklyanin-cubics.tex
% Current result: 2026-11-23 / v84.
\subsubsection{分離原点の Clifford augmentation に対する明示的な Shamash 模型}
\label{ifj:subsubsec:explicit-shamash}

混合係数が零のとき、原点の因子比較には
一因子の augmentation stabilization を使用する。
その存在だけでなく、指定中心三次の scalar と内部次数を保持した
有限行列を明示する。以下では右加群の行列を左乗法で記す。

\begin{proposition}[反転・巡回同変な Clifford 因子の四階 Shamash 行列分解]
\label{ifj:prop:clifford-explicit-shamash}
$\tau(\tau^3-1)(\tau^3+8)\ne0$ とし、
\[
 S_\tau=\C\langle x_1,x_2,x_3\rangle/
 (x_{i+1}x_{i+2}+x_{i+2}x_{i+1}+\tau x_i^2)_{i\in\mathbb Z/3}
\]
と置く（添字は巡回的に読む）。$\Delta=\tau^3-1$、
\[
 \omega=x_1x_2x_3+\frac{\tau}{2}(x_1^3-x_2^3+x_3^3),
 \qquad \Omega=2(1-\tau^3)\omega=-2\Delta\omega
\]
を指定中心三次とする。
$S_\tau$ の augmentation の三段 Koszul 分解を
\[
 P_\bullet:\quad
 0\longrightarrow S_\tau(-3)
 \xrightarrow{d_3}S_\tau(-2)^3
 \xrightarrow{d_2}S_\tau(-1)^3
 \xrightarrow{d_1}S_\tau\longrightarrow\C\longrightarrow0
\]
と書く。ただし
\[
 d_1=(x_1\ \ x_2\ \ x_3),\qquad
 d_2=\begin{pmatrix}
 \tau x_1&x_3&x_2\\x_3&\tau x_2&x_1\\x_2&x_1&\tau x_3
 \end{pmatrix},\qquad d_3=d_1^{\mathsf t}.
\]
次数二の列 $u=(u_i)_{i=1}^3$ と三階行列 $A$、
次数三の元 $v$ を、$\kappa=(4-\tau^3)/6$ を使って定める：
\begin{equation}\label{ifj:eq:explicit-shamash-data}
 \begin{aligned}
 u_i&=-\frac{\Delta}{3}
            (2x_{i+1}x_{i+2}+\tau x_i^2),\\
 A_{ii}&=-\frac{\tau^3}{2}x_i^2
                -\tau^2x_{i+1}x_{i+2},\\
 A_{i,i+1}&=\tau x_{i+2}^2+\kappa x_ix_{i+1},
 & A_{i+1,i}&=-\tau x_{i+2}^2-\kappa x_{i+1}x_i,\\
 v&=-\frac{\tau\Delta(\tau^3+8)}{18}
                         (x_1^3+x_2^3+x_3^3).
 \end{aligned}
\end{equation}
すると以下の恒等式が $S_\tau$ の中で厳密に成立する：
\begin{equation}\label{ifj:eq:explicit-shamash-identities}
 \begin{gathered}
 d_1u=u^{\mathsf t}d_3=\Omega,\\
 d_2A+ud_1=A d_2+d_3u^{\mathsf t}=\Omega\mathbf1_3,\\
 Au+d_3v=0,\qquad u^{\mathsf t}A+vd_1=0.
 \end{gathered}
\end{equation}
特に、$s_0=d$、$s_1=(u,A,u^{\mathsf t})$、$s_2=v$ は
指定中心元 $\Omega$ に対する higher homotopies である。

$F^0=S_\tau\oplus S_\tau(1)^3$、
$F^{-1}=S_\tau(-1)^3\oplus S_\tau$ と置き、
\begin{equation}\label{ifj:eq:explicit-shamash-mf}
 D_+=\begin{pmatrix}d_1&0\\A&d_3\end{pmatrix},
 \qquad
 D_-=\begin{pmatrix}u&d_2\\v&u^{\mathsf t}\end{pmatrix}
\end{equation}
とする。これらはそれぞれ次数零の射
$D_+:F^{-1}\to F^0$、$D_-:F^0\to F^{-1}(3)$ であり、
適切な twist を省略した行列の積は
\[
 D_+D_-=D_-D_+=\Omega\mathbf1_4
\]
となる。従ってこの $4\times4$ 行列分解は
$\operatorname{stab}_{S_\tau,\Omega}(\C)$ を表し、
指定した中心三次による内部三周期を保持する。
さらに巡回写像 $\rho:x_i\mapsto x_{i+1}$ と反転
\begin{equation}\label{ifj:eq:explicit-shamash-reflection}
 \sigma:(x_1,x_2,x_3)\longmapsto(-x_2,-x_1,-x_3)
\end{equation}
に対し、両微分と higher homotopies が厳密同変である。
$P$ を第一、第二座標の交換行列とすると、
$\sigma$ は Koszul 複体の四項に
\begin{equation}\label{ifj:eq:explicit-shamash-top-sign}
 (T_0,T_1,T_2,T_3)=(1,-P,P,-1)
\end{equation}
として作用し、特に最上 syzygy 線では $-1$ である。
$E_0=\operatorname{diag}(1,P)$、
$E_1=\operatorname{diag}(-P,-1)$ とすれば、
\begin{equation}\label{ifj:eq:explicit-shamash-equivariant-differentials}
 \sigma(D_+)=E_0D_+E_1^{-1},\qquad
 \sigma(D_-)=E_1D_-E_0^{-1}.
\end{equation}
$\rho^3=\sigma^2=1$、$\sigma\rho\sigma=\rho^{-1}$ も
これらの複体作用として厳密に成立する。
可換点曲線に作用する反転は $P$ であって $-P$ ではないが、
その最上 Koszul 線の符号 $\det P=-1$ と
式\eqref{ifj:eq:explicit-shamash-top-sign}の符号は一致する。
ここから比較核自体の同変性はまだ結論しない。

任意の中心 scalar $a$ に対して
$(u,A,v)\mapsto(au,aA,a^2v)$ とすれば、
曲率 $a\Omega$ の同じ模型となる。
従って分離二因子 $S_{\tau_1}\otimes S_{\tau_2}$ と
$w_0=a_1\Omega_{\tau_1}+a_2\Omega_{\tau_2}$ に対し、
二つのこの模型の Koszul 符号付き外積は
共通の三周期を同定した次数付き $32\mid32$ の
augmentation 行列分解を与える。
ただし、これを混合係数を動かした実 Ku 比較核と
同一視するものではない。
\end{proposition}

\begin{proof}
$S_\tau$ の滑らかな Sklyanin 領域では三段 Koszul 分解が完全であり、
$d_1d_2=d_2d_3=0$ は表示した三つの二次関係そのものである。
中心元 $\Omega$ の式と正則性は既存の Clifford 中心計算による。

式\eqref{ifj:eq:explicit-shamash-identities}は
$x_jx_i=-x_ix_j-\tau x_k^2$（$i<j$、$k$ は残る添字）および
中心平方 $x_i^2$ を使って順序 PBW 基底
\[
 1,x_1,x_2,x_3,x_1x_2,x_1x_3,x_2x_3,x_1x_2x_3
\]
へ整列して得る。
例えば巡回平均の恒等式
\[
 \omega=\frac13\sum_{i=1}^3
   \left(x_ix_{i+1}x_{i+2}+\frac{\tau}{2}x_i^3\right)
\]
から $d_1u=\Omega$ が直ちに従う。
同じ整列で $u^{\mathsf t}d_3=\Omega$、
$3\times3$ の二つの式および次数四の二つの式を確認する。
後二式は $s_1^2$ の nullhomotopy が $s_2=v$ であることを示す。
実際、$Au$ の第一成分の単位 PBW 成分に含まれる
$x_1^4$ の係数は
\[
 \frac{\tau(\tau^3-1)(\tau^3+8)}{18}\ne0
\]
なので、ここで選んだ巡回同変な $s_1$ では
$s_1^2$ を零とみなせない。
三段複体では $s_1s_2+s_2s_1=s_2^2=0$ が次数の理由で成立する。
従って主稿の Shamash folding の補題
\paperref{main}{mbnp:lem:shamash}を適用でき、
式\eqref{ifj:eq:explicit-shamash-mf}の二つの行列積と
augmentation stabilization の同定が従う。
各項の内部次数は $d_1,d_2,d_3$ が一次、
$u,A$ が二次、$v,\Omega$ が三次であり、
表示した $F^0,F^{-1}$ と整合する。
さらに $\sigma$ は関係 $r_i$ を $r_{P(i)}$ に送る一方で、
生成元 $x_i$ を $-x_{P(i)}$ に送る。従って
$\sigma(d_1)=d_1(-P)$、
$\sigma(d_2)=(-P)d_2P$、
$\sigma(d_3)=(-P)d_3$ であり、
式\eqref{ifj:eq:explicit-shamash-top-sign}は複体写像の条件から従う。
また式\eqref{ifj:eq:explicit-shamash-data}を整列すれば
\[
 \sigma(u)=(-P)u,\qquad
 \sigma(A)=P A(-P),\qquad \sigma(v)=-v
\]
を得る。ここで $x_ix_j$ と $x_jx_i$ を安易に交換せず、
二次関係を使っている。
これを折り畳んだ行列が
式\eqref{ifj:eq:explicit-shamash-equivariant-differentials}である。
巡回置換についても同様で、行列の積で
$\sigma\rho\sigma=\rho^{-1}$ を確認する。
最後の外積は両因子の奇次数写像の交差項を
Koszul 符号で打ち消す標準構成である。
\end{proof}

\begin{remark}[明示した範囲]
この行列分解は原点の\emph{非可換側}に存在する
実有限 Koszul--Shamash 模型である。
新しく構成したのは反転と巡回を同時に厳密に保つ $s_1,s_2$ と
その規格化された行列、最上 Koszul 線の符号であり、
一般的な Shamash folding は既知理論である。
分離因子の既存 window 同値における
$\operatorname{stab}(\C)\leftrightarrow\OO_E[1]$ には
主稿の系\paperref{main}{mbrn:cor:factor-boundary}を用いる。
この有限 MF 自体から、混合族の比較 kernel、
その Hom 複体の special retract $(i,p,h)$、
unit/counit、comparison の rotation と三乗 period の整合、
四つの probe の原式46枠と十二係数はまだ得られない。
有限多項式恒等式の再現は
\texttt{computations/verify\_clifford\_explicit\_shamash.py} に置く。
\end{remark}
